% GENERATED FILE. Edit canonical lessons, then run npm run generate:books.
% canonical-source-hash: fnv1a64:89dab0546a9185ca
% canonical-lessons: IT-C433-matrimoniale, IT-C433-culla, IT-C433-federa, IT-C433-saponetta, IT-C433-asciugacapelli

\chapter{In camera}
\label{ch:it-in-camera}
\hlchaptermodality{433}

\begin{chapteropening}
\textbf{By the end of this chapter:} \emph{I can ask for things in my room: a double room, a cot, a pillowcase, soap and a hair dryer.}

Complete the last lesson of chapter 433: I can ask for things in my room: a double room, a cot, a pillowcase, soap and a hair dryer.
\end{chapteropening}

\section[matrimoniale]{matrimoniale — double (of a bed or room); of marriage}

\label{lesson:IT-C433-matrimoniale}



\begin{quote}
Before the new one: say the Italian for to stay, then the Italian for the porter.
\end{quote}

\subsection*{You'll want to know: matrimoniale}
\textbf{matrimoniale} — \textquotedblleft{}double (of a bed or room); of marriage\textquotedblright{}.

\begin{quote}
Vorrei una camera matrimoniale per due notti.
\end{quote}

\begin{quote}
\emph{I would like a double room for two nights.}
\end{quote}

\emph{Un letto matrimoniale}; a room with two beds is \emph{una doppia}.

\begin{cousinweb}
\emph{matrimoniale} is built on \emph{matrimonio} 'marriage', a learned borrowing from Latin \emph{mātrimōnium}, built on \emph{māter} 'mother' (Lewis \& Short).
\end{cousinweb}

\begin{grammarlens}[title={What you've built}]
One more word for telling, explaining and describing at B1.
\end{grammarlens}

\subsection*{Your turn}
\begin{itemize}
\raggedright
  \item \emph{Say it:} \emph{matrimoniale}
  \item \emph{Say it:} \emph{matrimoniale}, once more
  \item \emph{Say it:} a sentence of your own with \emph{matrimoniale}
  \item \emph{Recall:} say the Italian for to stay, then the Italian for the porter, then say \emph{matrimoniale} again
\end{itemize}

\begin{tcolorbox}[breakable,colback=teal!4,colframe=teal!35!black,title={Before you move on}]
What does matrimoniale mean? (\textquotedblleft{}Double (of a bed or room); of marriage\textquotedblright{}.) Say it once more.
\end{tcolorbox}
\section[la culla]{la culla — the cot, the cradle}

\label{lesson:IT-C433-culla}



\begin{quote}
Before the new one: say the Italian for the porter, then the Italian for double.
\end{quote}

\subsection*{You'll want to know: la culla}
\textbf{la culla} — \textquotedblleft{}the cot, the cradle\textquotedblright{}.

\begin{quote}
Avete una culla per il bambino?
\end{quote}

\begin{quote}
\emph{Do you have a cot for the baby?}
\end{quote}

Figuratively: \emph{la culla della civiltà} — the cradle of civilisation.

\begin{cousinweb}
\emph{culla} is inherited from Vulgar Latin \emph{cunla} (EtymDB), a little form of Latin \emph{cūnae} 'cradle'.
\end{cousinweb}

\begin{grammarlens}[title={What you've built}]
One more word for telling, explaining and describing at B1.
\end{grammarlens}

\subsection*{Your turn}
\begin{itemize}
\raggedright
  \item \emph{Say it:} \emph{la culla}
  \item \emph{Say it:} \emph{la culla}, once more
  \item \emph{Say it:} a sentence of your own with \emph{la culla}
  \item \emph{Recall:} say the Italian for the porter, then the Italian for double, then say \emph{la culla} again
\end{itemize}

\begin{tcolorbox}[breakable,colback=teal!4,colframe=teal!35!black,title={Before you move on}]
What does la culla mean? (\textquotedblleft{}The cot, the cradle\textquotedblright{}.) Say it once more.
\end{tcolorbox}
\section[la federa]{la federa — the pillowcase}

\label{lesson:IT-C433-federa}



\begin{quote}
Before the new one: say the Italian for double, then the Italian for the cot.
\end{quote}

\subsection*{You'll want to know: la federa}
\textbf{la federa} — \textquotedblleft{}the pillowcase\textquotedblright{}.

\begin{quote}
La federa del cuscino è macchiata.
\end{quote}

\begin{quote}
\emph{The pillowcase is stained.}
\end{quote}

\emph{Il lenzuolo} is the sheet.

\begin{cousinweb}
\emph{federa} was borrowed from Lombardic \emph{fedara} 'feather' (EtymDB): first the cover of a feather pillow.
\end{cousinweb}

\begin{grammarlens}[title={What you've built}]
One more word for telling, explaining and describing at B1.
\end{grammarlens}

\subsection*{Your turn}
\begin{itemize}
\raggedright
  \item \emph{Say it:} \emph{la federa}
  \item \emph{Say it:} \emph{la federa}, once more
  \item \emph{Say it:} a sentence of your own with \emph{la federa}
  \item \emph{Recall:} say the Italian for double, then the Italian for the cot, then say \emph{la federa} again
\end{itemize}

\begin{tcolorbox}[breakable,colback=teal!4,colframe=teal!35!black,title={Before you move on}]
What does la federa mean? (\textquotedblleft{}The pillowcase\textquotedblright{}.) Say it once more.
\end{tcolorbox}
\section[la saponetta]{la saponetta — the bar of soap}

\label{lesson:IT-C433-saponetta}



\begin{quote}
Before the new one: say the Italian for the cot, then the Italian for the pillowcase.
\end{quote}

\subsection*{You'll want to know: la saponetta}
\textbf{la saponetta} — \textquotedblleft{}the bar of soap\textquotedblright{}.

\begin{quote}
In bagno manca la saponetta.
\end{quote}

\begin{quote}
\emph{There is no soap in the bathroom.}
\end{quote}

A little \emph{sapone}.

\begin{cousinweb}
\emph{saponetta} is a little \emph{sapone} 'soap', which is inherited from Latin \emph{sāpō, sāpōnis} (EtymDB), itself an early loan from a Germanic language (Lewis \& Short).
\end{cousinweb}

\begin{grammarlens}[title={What you've built}]
One more word for telling, explaining and describing at B1.
\end{grammarlens}

\subsection*{Your turn}
\begin{itemize}
\raggedright
  \item \emph{Say it:} \emph{la saponetta}
  \item \emph{Say it:} \emph{la saponetta}, once more
  \item \emph{Say it:} a sentence of your own with \emph{la saponetta}
  \item \emph{Recall:} say the Italian for the cot, then the Italian for the pillowcase, then say \emph{la saponetta} again
\end{itemize}

\begin{tcolorbox}[breakable,colback=teal!4,colframe=teal!35!black,title={Before you move on}]
What does la saponetta mean? (\textquotedblleft{}The bar of soap\textquotedblright{}.) Say it once more.
\end{tcolorbox}
\section[l'asciugacapelli]{l'asciugacapelli — the hair dryer}

\label{lesson:IT-C433-asciugacapelli}



\begin{quote}
Before the new one: say the Italian for the pillowcase, then the Italian for the bar of soap.
\end{quote}

\subsection*{You'll want to know: l'asciugacapelli}
\textbf{l'asciugacapelli} — \textquotedblleft{}the hair dryer\textquotedblright{}.

\begin{quote}
L'asciugacapelli è nel cassetto.
\end{quote}

\begin{quote}
\emph{The hair dryer is in the drawer.}
\end{quote}

Plural unchanged; in speech also \emph{il phon}.

\begin{cousinweb}
\emph{asciugacapelli} is built from \emph{asciugare} 'to dry', inherited from Latin \emph{exsūcāre} 'to suck the juice out', and \emph{capelli} 'hair', inherited from Latin \emph{capillus} (EtymDB).
\end{cousinweb}

\begin{grammarlens}[title={What you've built}]
One more word for telling, explaining and describing at B1.
\end{grammarlens}

\subsection*{Your turn}
\begin{itemize}
\raggedright
  \item \emph{Say it:} \emph{l'asciugacapelli}
  \item \emph{Say it:} \emph{l'asciugacapelli}, once more
  \item \emph{Say it:} a sentence of your own with \emph{l'asciugacapelli}
  \item \emph{Recall:} say the Italian for the pillowcase, then the Italian for the bar of soap, then say \emph{l'asciugacapelli} again
\end{itemize}

\begin{tcolorbox}[breakable,colback=teal!4,colframe=teal!35!black,title={Before you move on}]
What does l'asciugacapelli mean? (\textquotedblleft{}The hair dryer\textquotedblright{}.) Say it once more.
\end{tcolorbox}
